\documentclass[conference]{IEEEtran}
\usepackage{float}
\usepackage{cite}
\usepackage{multirow}
\usepackage{amsmath}
\usepackage{makecell}
\usepackage{url}
\usepackage{graphicx}
\usepackage[hidelinks,bookmarks=false]{hyperref}
\usepackage[utf8]{inputenc}
\usepackage[spanish,es-noshorthands]{babel}
\usepackage{textcomp}

\ifCLASSINFOpdf
\else
\fi

\begin{document}

\markboth{ITCR. Escuela de Ingeniería Electrónica}{Tracker LoRa-APRS para monitoreo de bicicletas}

\title{Informe Final: Tracker LoRa-APRS para Monitoreo de Bicicletas de Montaña}

\author{\IEEEauthorblockN{1\textsuperscript{st} Fabián Villegas Bonilla}
\IEEEauthorblockA{\textit{Taller Integrador} \\
\textit{Instituto Tecnológico de Costa Rica}\\
Cartago, Costa Rica \\
2018235959 \\
famavibo@estudiantec.cr}
\and
\IEEEauthorblockN{2\textsuperscript{nd} Jesús Soto Benavides}
\IEEEauthorblockA{\textit{Taller Integrador} \\
\textit{Instituto Tecnológico de Costa Rica}\\
Cartago, Costa Rica \\
jjsoto1222@estudiantec.cr}}

\maketitle

% ============================================================
% RESUMEN
% ============================================================
\begin{abstract}
Este informe presenta el desarrollo de un tracker GPS autónomo basado en la plataforma T-Beam V1.2 (ESP32, LoRa SX1276, GPS NEO-6M) para el monitoreo de bicicletas de montaña en zonas sin cobertura celular. El sistema utiliza el protocolo APRS (Automatic Packet Reporting System) sobre LoRa para transmitir la posición en tiempo real a una frecuencia de 433.775 MHz. Se implementó un firmware optimizado para bajo consumo energético, con ciclos de transmisión de 60 segundos y deep sleep, logrando una autonomía estimada de 25 a 32 horas con batería de 1200 mAh. El proyecto incluye el diseño de hardware, desarrollo de firmware, pruebas de campo y análisis de costos para su potencial comercialización en grupos ciclistas.

\vspace{0.5cm}
\begin{IEEEkeywords}
LoRa, APRS, tracker GPS, ESP32, T-Beam, ciclismo de montaña, bajo consumo, radioaficionado
\end{IEEEkeywords}
\end{abstract}

% ============================================================
% INTRODUCCIÓN
% ============================================================
\section{Introducción}

El ciclismo de montaña es una actividad que se desarrolla frecuentemente en zonas rurales y montañosas con nula o limitada cobertura de telefonía celular. Esta situación representa un riesgo para los ciclistas, ya que en caso de accidente o emergencia no es posible comunicarse con servicios de rescate o compartir la ubicación con otros miembros del grupo.

El sistema APRS (Automatic Packet Reporting System) es un protocolo de comunicación diseñado originalmente para radioaficionados que permite transmitir información de posición, meteorológica y mensajes cortos a través de radiofrecuencia. Tradicionalmente opera en bandas VHF (144.390 MHz en Norteamérica), pero recientemente se ha adaptado para funcionar sobre LoRa (Long Range), una tecnología de modulación que ofrece gran alcance con bajo consumo energético.

El presente proyecto desarrolla un tracker GPS autónomo basado en la placa T-Beam V1.2, que integra un microcontrolador ESP32, un módulo LoRa SX1276, un GPS NEO-6M y un gestor de energía AXP2101. El firmware implementa el protocolo APRS sobre LoRa en la frecuencia de 433.775 MHz (banda de 70 cm para radioaficionados en Costa Rica), permitiendo que la posición del ciclista sea visible en plataformas como APRS.fi a través de estaciones iGate cuando se encuentren dentro del rango de cobertura.

Este informe detalla la arquitectura del sistema, el desarrollo del firmware, las pruebas de campo realizadas, los resultados obtenidos y un análisis de costos para la producción de 6 unidades dirigidas a grupos ciclistas.

% ============================================================
% OBJETIVOS
% ============================================================
\section{Objetivos}

\subsection{Objetivo General}
Desarrollar un tracker GPS de bajo consumo basado en tecnología LoRa-APRS para el monitoreo de bicicletas de montaña en zonas sin cobertura celular.

\subsection{Objetivos Específicos}
\begin{enumerate}
    \item Implementar un firmware para la placa T-Beam V1.2 que obtenga posición GPS y la transmita en formato APRS sobre LoRa.
    \item Optimizar el consumo energético mediante deep sleep del ESP32 y gestión eficiente de periféricos.
    \item Validar el funcionamiento del tracker mediante pruebas de campo en entornos reales de ciclismo.
    \item Documentar el diseño y desarrollo para su reproducción por parte de la comunidad de radioaficionados y ciclistas.
\end{enumerate}

% ============================================================
% ARQUITECTURA DEL SISTEMA
% ============================================================
\section{Arquitectura del Sistema}

\subsection{Plataforma Hardware}

La placa T-Beam V1.2 integra todos los componentes necesarios en un formato compacto:

\begin{table}[H]
    \centering
    \caption{Componentes principales del T-Beam V1.2}
    \begin{tabular}{|c|c|c|}
        \hline
        \textbf{Componente} & \textbf{Modelo} & \textbf{Función} \\
        \hline
        Microcontrolador & ESP32 & Procesador principal y gestión de energía \\
        LoRa & SX1276 & Transceptor de 433/868/915 MHz \\
        GPS & NEO-6M & Receptor satelital con precisión \textless 2.5m \\
        PMU & AXP2101 & Gestión de energía y carga de batería \\
        OLED & SSD1306 & Pantalla de 128x64 píxeles \\
        Batería & LiPo 1000mAh & Alimentación autónoma \\
        \hline
    \end{tabular}
    
    \label{tab:hardware}
\end{table}

\subsection{Software y Protocolos}

\begin{itemize}
    \item \textbf{PlatformIO}: Entorno de desarrollo sobre VS Code
    \item \textbf{TinyGPSPlus}: Librería para parseo de tramas NMEA del GPS
    \item \textbf{LoRa.h}: Librería de bajo nivel para el SX1276
    \item \textbf{XPowersLib}: Control del AXP2101 (energía)
    \item \textbf{Adafruit GFX/SSD1306}: Manejo de la pantalla OLED
    \item \textbf{Protocolo APRS}: Formato estándar para la transmisión de posición
\end{itemize}

% ============================================================
% DESARROLLO DEL FIRMWARE
% ============================================================
\section{Desarrollo del Firmware}

\subsection{Características Principales}

El firmware desarrollado consta de aproximadamente 500 líneas de código en C++, organizado en las siguientes rutinas principales:

\begin{enumerate}
    \item \textbf{Inicialización secuencial}: AXP2101 - OLED - Botón - GPS - LoRa
    \item \textbf{Espera de fix GPS}: Timeout de 60 segundos con visualización de satélites
    \item \textbf{Lectura de datos}: Extracción de latitud, longitud, velocidad, rumbo y altitud
    \item \textbf{Conversión a formato APRS}: Coordenadas DDMM.mm y altitud en pies
    \item \textbf{Transmisión LoRa}: 3 reintentos automáticos con preámbulo \textless 0xFF 0x01
    \item \textbf{Deep sleep}: 60 segundos con wakeup por timer o botón
\end{enumerate}

\subsection{Eficiencia Energética}

Se implementaron múltiples estrategias para maximizar la duración de la batería:

\begin{table}[H]
    \centering
    \caption{Optimizaciones de consumo energético}
    \begin{tabular}{|c|c|c|}
        \hline
        \textbf{Estrategia} & \textbf{Configuración} & \textbf{Ahorro estimado} \\
        \hline
        Frecuencia CPU & 80 MHz (vs 240 MHz) & 40\% en modo activo \\
        Deep sleep & 60 segundos entre ciclos & 95\% del tiempo apagado \\
        OLED apagada & En sleep y al finalizar & 99\% menos consumo \\
        GPS persistente & Permanece encendido & Fix en \textless 5s \\
        \hline
    \end{tabular}
    
    \label{tab:energia}
\end{table}

\subsection{Máquina de Estados}

El flujo principal del firmware se representa mediante la siguiente máquina de estados:

\begin{enumerate}
    \item \textbf{Inicio}: Encendido o wakeup (timer o botón)
    \item \textbf{Inicialización}: Configuración de periféricos
    \item \textbf{Espera GPS}: Búsqueda de fix satelital (máx 60s)
    \item \textbf{Transmisión}: Envío de trama APRS por LoRa (3 intentos)
    \item \textbf{Datos finales}: Visualización de posición (5s si TX exitoso)
    \item \textbf{Deep sleep}: Suspensión por 60 segundos
    \item \textbf{Retorno al paso 1}: Wakeup por timer o botón
\end{enumerate}

% ============================================================
% PRUEBAS Y RESULTADOS
% ============================================================
\section{Pruebas y Resultados}

\subsection{Prueba del firmware propio}

Se realizaron pruebas del firmware desarrollado para validar la correcta inicialización de los periféricos, la obtención del fix GPS y la transmisión LoRa. La salida del monitor serial, que se muestra en el Apéndice A (\hyperlink{apendice:serial}{ver}), evidencia los siguientes resultados:

\begin{itemize}
    \item Inicialización correcta del gestor de energía AXP2101 con voltaje de batería de 4.13 V y 100\% de carga.
    \item Fix GPS obtenido en aproximadamente 2.3 segundos con 8 satélites visibles.
    \item Transmisión LoRa exitosa, identificada con el mensaje \texttt{TX OK — transmision \#22}.
    \item Construcción correcta de la trama APRS con el formato estándar.
\end{itemize}

\subsection{Demostración del funcionamiento en pantalla}

El video presentado en el Apéndice B (\hyperlink{apendice:video}{ver}) muestra el comportamiento del tracker en tiempo real. En él se observa:

\begin{enumerate}
    \item El mensaje de bienvenida en la pantalla OLED al encender el dispositivo.
    \item La obtención del fix GPS con visualización de coordenadas (latitud y longitud).
    \item La alternancia automática entre la pantalla de posición y la pantalla de estado del sistema cada 3 segundos.
    \item La indicación del estado de transmisión LoRa en la pantalla de estado.
    \item La entrada en deep sleep al finalizar el ciclo completo.
\end{enumerate}

Para confirmar que las tramas transmitidas mediante LoRa llegaban correctamente a internet, se verificó la plataforma APRS.fi. Como se muestra en el Apéndice C (\hyperlink{apendice:aprs}{ver}), el tracker con callsign “TI0TEC7-7” apareció registrado en el mapa junto con su posición geográfica, lo que confirma que la trama APRS generada por el firmware posee el formato adecuado.

Además, la prueba permitió comprobar la existencia de al menos un iGate (Internet Gateway) dentro del rango de cobertura capaz de recibir la transmisión y reenviarla a la red APRS-IS. Asimismo, se validó que la frecuencia de 433.775 MHz y los parámetros LoRa utilizados (SF12, BW125kHz y CR4/5) son compatibles con la infraestructura APRS disponible en Costa Rica.

\subsection{Máquina de estados del sistema}

El funcionamiento del firmware se organiza en una máquina de estados, representada gráficamente en el Apéndice D (\hyperlink{apendice:maquina}{ver}). Los estados principales y sus transiciones son:

\begin{table}[H]
    \centering
    \scriptsize
    \caption{Descripción de los estados del sistema}
    \begin{tabular}{|c|l|p{4cm}|}
        \hline
        \textbf{Estado} & \textbf{Nombre} & \textbf{Descripción} \\
        \hline
        1 & Inicialización & Configuración secuencial de periféricos (AXP2101 - OLED - Botón - GPS - LoRa) \\
        \hline
        2 & Espera GPS & Búsqueda de fix satelital con timeout de 60 segundos \\
        \hline
        3 & Lectura GPS & Extracción de latitud, longitud, velocidad, rumbo y altitud \\
        \hline
        4 & Construcción APRS & Conversión de coordenadas a formato DDMM.mm y ensamblaje de trama \\
        \hline
        5 & Transmisión LoRa & Envío de la trama con hasta 3 reintentos automáticos \\
        \hline
        6 & Datos finales & Visualización de posición en pantalla (5 segundos, solo si TX fue exitoso) \\
        \hline
        7 & Deep Sleep & Suspensión del ESP32 por 60 segundos con wakeup por timer o botón \\
        \hline
    \end{tabular}
    \label{tab:estados}
\end{table}

Las transiciones más relevantes son:
\begin{itemize}
    \item \textbf{Fix válido}: Del estado 2 (Espera GPS) al estado 3 (Lectura GPS)
    \item \textbf{Timeout sin fix}: Del estado 2 directamente al estado 7 (Deep Sleep)
    \item \textbf{TX exitoso}: Del estado 5 (Transmisión LoRa) al estado 6 (Datos finales)
    \item \textbf{TX fallido}: Del estado 5 al estado 7 después de 3 reintentos
    \item \textbf{Wakeup por timer}: Cada 60 segundos, del estado 7 al estado 1
    \item \textbf{Wakeup por botón}: Presión del GPIO38, del estado 7 al estado 1
\end{itemize}

Este diseño garantiza que el tracker siempre complete un ciclo completo y nunca se quede atascado en un estado, maximizando la eficiencia energética.

\subsection{Resumen de métricas obtenidas}

\begin{table}[H]
    \centering
    \scriptsize
    \caption{Resultados de las pruebas de campo}
    \begin{tabular}{|c|c|c|}
        \hline
        \textbf{Métrica} & \textbf{Resultado} & \textbf{Condición} \\
        \hline
        Tiempo de fix GPS & 2.3 segundos & Exterior, cielo despejado \\
        Satélites visibles & 8 & Exterior \\
        Tasa de éxito TX & $>$ 90\% & Zona urbana (1-2 km) \\
        Autonomía estimada & 25-36 horas & Batería 1200 mAh \\
        \hline
    \end{tabular}
    \label{tab:resultados}
\end{table}

\subsection{Soporte bicicleta}

Se realizo un diseño 3D para introducir el tracker un soporte para bicis visto de (\hyperlink{apendice:soporte}{ver}) el cual tiene la función llevar el tracker en todo momento en la bicicleta para no molestar con peso extra al usuario de esta y que este firme y segura de que nada le pueda pasar.


% ============================================================
% ESTIMACIÓN DE COSTOS
% ============================================================
\section{Estimación de Costos}

\subsection{Hardware}

Los componentes físicos son adquiridos por el equipo desarrollador e incluidos en el precio final del producto entregado a cada grupo de ciclistas.

\begin{table}[H]
    \centering
    \scriptsize
    \caption{Costos de hardware para 6 unidades}
    \begin{tabular}{|c|c|c|c|}
        \hline
        \textbf{Componente} & \textbf{Cantidad} & \textbf{P. Unitario(₡)} & \textbf{Total(₡)} \\
        \hline
        Tracker T-Beam V1.2 & 6 & 31.000 & 186.000 \\
        Soporte de montaje para bicicleta & 6 & 10.000 & 60.000 \\
        \hline
        \textbf{Subtotal hardware} & & & \textbf{246.000} \\
        \hline
    \end{tabular}
    \label{tab:costo_hardware}
\end{table}

\subsection{Honorarios de Ingeniería}

Tarifa oficial ₡37.700 por hora profesional (CFIA, La Gaceta N°224, revisada enero 2026).

\begin{table}[H]
    \centering
    \scriptsize
    \caption{Honorarios de ingeniería según tarifa CFIA}
    \begin{tabular}{|l|c|c|r|}
        \hline
        \textbf{Fase / Actividad} & \textbf{Horas} & \textbf{Tarifa/hr(₡)} & \textbf{Subtotal(₡)} \\
        \hline
        Análisis de requerimientos & 5 & 37.700 & 188.500 \\
        Diseño del firmware y arquitectura APRS & 8 & 37.700 & 301.600 \\
        Desarrollo del firmware & 20 & 37.700 & 754.000 \\
        Configuración APRS y validación & 6 & 37.700 & 226.200 \\
        Pruebas de campo y documentación & 10 & 37.700 & 377.000 \\
        \hline
        \multicolumn{2}{|c|}{\textbf{Total}} & & \textbf{1.847.300} \\
        \hline
    \end{tabular}
    \label{tab:honorarios}
\end{table}

\subsection{Resumen Financiero}

\begin{table}[H]
    \centering
    \caption{Resumen financiero total del proyecto}
    \begin{tabular}{|c|c|}
        \hline
        \textbf{Concepto} & \textbf{Monto(₡)} \\
        \hline
        Hardware (trackers + soportes × 6) & 246.000 \\
        Honorarios de ingeniería (49 horas) & 1.847.300 \\
        Backend / Frontend (APRS.fi) & 0 \\
        \hline
        Subtotal sin IVA & 2.093.300 \\
        IVA 13\% (sobre honorarios) & 240.149 \\
        \hline
        \textbf{Total del Proyecto} & \textbf{2.333.449} \\
        \hline
    \end{tabular}
    
    \label{tab:resumen}
\end{table}

El IVA del 13\% aplica únicamente sobre los honorarios profesionales conforme a la Ley N°9635. El uso de APRS.fi elimina el desarrollo de backend y frontend, reduciendo significativamente el costo total del proyecto.

% ============================================================
% CONCLUSIONES
% ============================================================
\section{Conclusiones}

El desarrollo del tracker LoRa-APRS para monitoreo de bicicletas de montaña ha cumplido con los objetivos planteados:

\begin{enumerate}
    \item Se implementó un firmware funcional en PlatformIO que obtiene posición GPS, construye tramas APRS válidas y las transmite por LoRa a 433.775 MHz.
    
    \item Se optimizó el consumo energético mediante CPU a 80 MHz, deep sleep de 60 segundos y apagado automático de la pantalla OLED, logrando una autonomía de 12-15 horas con batería de 1000 mAh.
    
    \item Las pruebas de campo demostraron un tiempo de fix GPS entre 2 y 30 segundos en exteriores, con una tasa de éxito de transmisión superior al 90\% en zonas urbanas con cobertura de iGates.
    
    \item Se documentó completamente el código, diagramas y procedimientos para permitir la reproducción del proyecto por terceros.
\end{enumerate}

\subsection{Trabajo Futuro}

Como mejoras potenciales se identifican:
\begin{itemize}
    \item Implementar recepción continua (RX) con interrupciones para detectar iGates cercanos
    \item Agregar almacenamiento de posición en memoria flash cuando no hay cobertura
    \item Desarrollar una aplicación móvil para configuración del tracker vía Bluetooth
    \item Integrar sensores adicionales (temperatura, ritmo cardíaco, cadencia)
\end{itemize}

% ============================================================
% REFERENCIAS
% ============================================================


\clearpage
\appendix
\renewcommand{\thesection}{\Alph{section}}
\setcounter{section}{0}

\section*{Apéndice A: Registro serial del firmware propio}
\label{apendice:serial}
\hypertarget{apendice:serial}{}

\begin{figure}[H]
    \centering
    \includegraphics[width=1\columnwidth]{prueba_firmware_propio_v1.jpeg}
    \caption{Salida del monitor serial durante la ejecución del firmware propio}
    \label{fig:serial}
\end{figure}


\section*{Apéndice B: Video demostrativo del funcionamiento}
\label{apendice:video}
\hypertarget{apendice:video}{}

\textbf{Enlace directo:} \url{https://www.youtube.com/shorts/KhTmpQa0nLU}

\section*{Apéndice C: Visualización en APRS.fi}
\label{apendice:aprs}
\hypertarget{apendice:aprs}{}

\begin{figure}[H]
    \centering
    \includegraphics[width=1\columnwidth]{aprs.png}
    \caption{Posición del tracker registrada en APRS.fi}
    \label{fig:aprs}
\end{figure}


\section*{Apéndice D: Máquina de estados del sistema}
\label{apendice:maquina}
\hypertarget{apendice:maquina}{}

\begin{figure}[H]
    \centering
    \includegraphics[width=0.8\columnwidth]{maquina_de_estados.png}
    \caption{Diagrama de la máquina de estados del firmware}
    \label{fig:maquina}
\end{figure}

\section*{Apéndice E: Diseño soporte}
\label{apendice:soporte}
\hypertarget{apendice:soporte}{}

\begin{figure}[H]
    \centering
    \includegraphics[width=0.8\columnwidth]{Integrador_1.PNG}
    \caption{Diagrama soporte isometrico}
    \label{fig:Soporte1}
\end{figure}

\begin{figure}[H]
    \centering
    \includegraphics[width=0.8\columnwidth]{Integrador_2.PNG}
    \caption{Diagrama soporte de frente}
    \label{fig:Soporte2}
\end{figure}

\begin{figure}[H]
    \centering
    
\includegraphics[width=0.8\columnwidth]{Integrador_3.PNG}
    \caption{Diagrama soporte de lado}
    \label{fig:Soporte3}
\end{figure}

\end{document}
